\documentclass[11pt]{article}
\usepackage[T1]{fontenc}
\usepackage{textcomp}
\usepackage[margin=0.75in]{geometry}
\usepackage{amsmath,amssymb,amsthm}
\usepackage{array,booktabs}
\usepackage{hyperref}
\setlength{\parindent}{0pt}
\newtheorem{theorem}{Theorem}
\theoremstyle{definition}
\newtheorem{definition}{Definition}
\title{Flow Proof Artifact: prop-03-rectangle-whole-and-other-part}
\author{Generated by \texttt{flow doc proof}}
\date{Flow Proof Book}
\begin{document}
\maketitle
If a straight line is cut, the rectangle by the whole and the other segment equals the rectangle by the segments together with the square on the other segment.\\[0.5em]
\textbf{Source.} euclid — Elements, Book II, Proposition 3\\[0.5em]
\bigskip
\noindent{\large \textbf{Derived fact 1.} \textit{If a straight line is cut, the rectangle by the whole and the other segment equals the rectangle by the segments together with the square on the other segment} \hfill the Euclidean plane \cdot Euclid Book II \cdot Proposition 3: the rectangle by the whole and the other part equals the parts plus a square}\par
\smallskip\noindent\textit{Coordinate: the Euclidean plane \cdot Euclid Book II \cdot Proposition 3: the rectangle by the whole and the other part equals the parts plus a square}\par
\smallskip\noindent\textit{Source: euclid — Elements, Book II, Proposition 3}\par
\smallskip\noindent\textit{Built on: proposition 2: the rectangle by the whole and one part equals the parts plus a square, for Euclid Book II on the Euclidean plane, proposition 34: complements of parallelograms about the diameter are equal, for Book I of the Elements in on the Euclidean plane}\par
\medskip
\noindent\textbf{Goal.} If a straight line is cut, the rectangle by the whole and the other segment equals the rectangle by the segments together with the square on the other segment\par
\noindent$\displaystyle rect(whole, \text{part D}) = rect(parts) + sq(D)$\par
\medskip
\noindent
\begin{tabular}{@{} >{\raggedright\arraybackslash}p{0.06\textwidth} >{\raggedright\arraybackslash}p{0.47\textwidth} >{\raggedleft\arraybackslash}p{0.41\textwidth} @{}}
\textbf{\#} & \textbf{Proof} & \textbf{Mathematics} \\
\midrule
\textbf{1.} & We prove that proposition 3: the rectangle by the whole and the other part equals the parts plus a square for Euclid Book II the Euclidean plane. & \\[0.45em]
\textbf{2.} & Let whole = straight\_line\_AB. & \\[0.45em]
\textbf{3.} & Let C = segment\_AC. & \\[0.45em]
\textbf{4.} & Let D = segment\_CB. & \\[0.45em]
\textbf{5.} & From step 2, step 3, and step 4, we can deduce that rect(whole, D) equals rect(C, D) plus sq(D). & $\displaystyle rect(whole, D) = rect(C, D) + sq(D)$ \\[0.45em]
\textbf{6.} & We invoke the derived fact governing Euclid Book II the Euclidean plane: proposition 2: the rectangle by the whole and one part equals the parts plus a square, for Euclid Book II on the Euclidean plane. & \\[0.45em]
\textbf{7.} & We invoke the derived fact governing Book I of the Elements in the Euclidean plane: proposition 34: complements of parallelograms about the diameter are equal, for Book I of the Elements in on the Euclidean plane. & \\[0.45em]
\textbf{8.} & From step 6 and step 7, this implies rect(whole, part D) equals rect(parts) plus sq(D). Hence proven. & $\displaystyle rect(whole, \text{part D}) = rect(parts) + sq(D)$ \\[0.45em]
\bottomrule
\end{tabular}
\label{thm:Geometry-Euclid Book II-proposition_3_the_rectangle_by_the_whole_and_the_other_part_equals_the_parts_plus_a_square}
\par\medskip\hrule
\end{document}
